\subsection{Técnica de Contribución por Elemento}

\graybold{Preguntas guía:}

\begin{itemize}
\item ¿Piden una suma agregada sobre TODAS las subsecuencias/subconjuntos posibles (potencialmente \bigO{n^2} o más de ellas)?
\item ¿Puedo, en vez de sumar sobre subsecuencias, sumar sobre ELEMENTOS, contando cuántas subsecuencias incluyen a cada uno?
\item ¿El "peso" (cantidad de subsecuencias que incluyen a un elemento) tiene una fórmula cerrada simple?
\end{itemize}

\graybold{¿Para qué sirve?} Convertir una suma sobre un número cuadrático (o mayor) de subconjuntos/subsecuencias en una suma lineal sobre los elementos originales, calculando cuántas veces (con qué peso) cada elemento contribuye al resultado final. Es una de las técnicas más rentables en matemáticas y simulación: cambia la perspectiva de "qué subconjuntos sumo" a "cuánto aporta cada elemento".

\graybold{Complejidad:} \bigO{n}.

\graybold{Fundamento teórico:} La suma de $S(L,R)$ sobre todos los pares $1 \le L \le R \le n$ se puede reescribir intercambiando el orden de sumatoria: en vez de sumar por subperíodo, se suma por elemento $a_i$, contando en cuántos subperíodos $[L,R]$ aparece. Un elemento en la posición $i$ (1-indexada) aparece en un subperíodo $[L,R]$ si y solo si $L \le i \le R$; hay exactamente $i$ opciones para $L$ (de 1 a $i$) y $(n-i+1)$ opciones para $R$ (de $i$ a $n$). Por lo tanto, el peso de $a_i$ es $i \cdot (n-i+1)$, y la respuesta final es $\sum_{i=1}^{n} a_i \cdot i \cdot (n-i+1)$.

\graybold{Ejercicio aplicado}

Dada una secuencia de n números, calcular la suma de $S(L,R)$ (la suma del subperíodo $[L,R]$) sobre TODOS los pares posibles $(L,R)$ con $1 \le L \le R \le n$.

\graybold{I/O:} n ≤ 10\textsuperscript{5} → lectura total con tokenización (patrón 2).

\graybold{Código solución}

\begin{lstlisting}[language=Python]
import sys

def solve():
    data = sys.stdin.buffer.read().split()
    n = int(data[0])
    a = list(map(int, data[1:1 + n]))
    total = 0
    for i in range(1, n + 1):
        total += a[i - 1] * i * (n - i + 1)   # peso = # subperiodos que incluyen la posicion i
    print(total)

solve()
\end{lstlisting}